\scanpage{80}%
valente à $\handwritten{\underline}{G}(\handwritten{\unsure{\ell|\gamma}})$ pour tout corps alg. clos $\handwritten{\ell}$ de car\deleted{\unsure{r|.}}. $p\handwritten{>0}$ (de façon précise, si $\handwritten{\overline}{\unsure{\handwritten{\underline}{F}|\handwritten{\underline}{\mathbb{F}}}}_p$ est\deleted{\unsure{x|t}} la clôture algébrique de $\handwritten{\underline}{F}_p$ dans $\handwritten{\ell}$, l'inclusion de celle-ci dans $\handwritten{\ell}$ induit une équivalence de gerbes
\[
\handwritten{\underline}{G}(\handwritten{\ell})\handwritten{\underline{\approx}}\handwritten{\underline}{G}(\handwritten{\overline}{\handwritten{\underline}{F}}_p)\deleted{\handwritten{\unsure{)|,}}}\handwritten{\underline{\approx}}\handwritten{\underline}{H}\deleted{\handwritten{\unsure{)|.}}}\qquad ).
\]
8.6. Etude de $\operatorname{Rep}(\handwritten{\underline}{H})$. Nous avons dit que les classes d'objets simples sont en correspondance biunivoque avec les éléments de $\handwritten{\underline}{Q}$ ; l'élément de $\handwritten{\underline}{Q}$ correspondant à un objet isotypique est dit la \handwritten{\uline}{pente} de cet élément. Nous avons vu que l'objet de $\operatorname{Rep}(\handwritten{\underline}{H})$ défini \margin{\handwritten{Notation $E_{s/r}$ pour $E_{r,s}$ indiquée, donc
\[
E_\lambda\otimes E_{\lambda'}\simeq E_{\lambda\lambda'}^{n(\lambda,\lambda')}
\]
pour $\lambda,\lambda'\in\mathbb{Q}$}}
par l'objet $\handwritten{\underline}{Q}_p[T]/(T^r-p^s)$ de $\operatorname{Fcriso}(\handwritten{\underline}{F}_p)$, muni de la multiplication par $T$, est simple de pente $s/r$ ($s,r\handwritten{\in}\handwritten{\underline}{Z}$, $r\geq1$, $(s,r)=1$ ), de sorte que nous obtenons ainsi un catalogue complet des objets simples $E_{r,s}$. Bien entendu, $E_{r,s}$ est de rang $r$, d'autre part son corps des endomorphismes est \deleted{\unsure{il|ila}} un corps gauche \inserted{$K_{r,s}$} dont la classe dans $\operatorname{Br}(\handwritten{\underline}{Q}_p)=\handwritten{\underline}{Q}/\handwritten{\underline}{Z}$ est $s/r\bmod1$, d'après ce qu'on a vu dans 8.3.

$\operatorname{Rep}(\handwritten{\underline}{H})$ est limite inductive des sous-catégories pleines $\operatorname{Rep}(\handwritten{\underline}{H},r)$, $r\geq1$, où $(\handwritten{\underline}{H},r)$ est le quotient de la gerbe $\handwritten{\underline}{H}$ correspond au quotient $(H,r)\handwritten{\simeq}\handwritten{\underline}{G}_m$ de son lien, tel que l'homomorphisme de Tate \deleted{\unsure{soit un|soitxxx}} induise un homomorphisme $(H,r)\handwritten{\xrightarrow{\unsure{a|q|\alpha}}}\handwritten{\underline}{G}_m$ s'identifiant à $x\handwritten{\longmapsto}x^r$ : F-iso\inserted{cristaux (sur $\handwritten{\overline}{k}$) de pente multiple de $1/r$.} Donc $(\handwritten{\underline}{H},r)$ \deleted{\unsure{s'identifie}} èst\deleted{\unsure{x|à}} \inserted{une} gerbe de lien $\handwritten{\underline}{G}_m$, ayant comme classe dans $\operatorname{Br}(\handwritten{\underline}{Q}_p)$ le générateur canonique $1/r\bmod\handwritten{\underline}{Z}$ de ${}_r\operatorname{Br}(\handwritten{\underline}{Q}_p)\handwritten{\simeq}\handwritten{\underline}{Z}/r\handwritten{\underline}{Z}$. C\unsure{,'|'}est aussi, plus canoniquement, ``la'' gerbe définie par le corps gauche $K_{r,\handwritten{1}\deleted{\unsure{s|1}}}$ plus haut, \deleted{qui}

\handwritten{[}Exercice amusant, que je n'ai pas fait: la structure additive de $\operatorname{Rep}(\handwritten{\underline}{H})$ étant entièrement décrite par la collection des $K_{r,s}$, expliciter sa $\handwritten{\otimes}$-structure ... i.e. écrire les \handwritten{``}relations canoniques\handwritten{''} en termes de prod\unsure{u|n}i\unsure{t|e}s tensoriels entre les $E_{r,s}$. Autre exerciee suggéré: trouver ``canoniquement'' une description des modules de Dieudonné des réduits des groupes de Lubin-Tate en termes des $E_{r,s}$ \deleted{\handwritten{\unsure{,|.}}} \handwritten{(}en particulier, les Lubin-Tate associés aux extensions non ramifiées
% archive sheet number: 79
